\documentclass[../cours_3e.tex]{subfiles}
\begin{document}
\setcounter{chapter}{19}
\chapter{Outils numériques et synthèse}
\label{ch:outils}
\begin{objectif}
Utiliser un tableur, une calculatrice et un logiciel de géométrie dynamique pour chercher, représenter et vérifier. Réunir plusieurs chapitres dans une même solution.
\end{objectif}
\section{Calculatrice : outil de calcul, pas de raisonnement}
Avant de saisir, estimer le résultat et écrire l'expression. Vérifier les parenthèses, les signes, l'unité d'angle et le rang d'arrondi. Conserver une valeur exacte lorsque c'est possible.
\begin{exemple}
Pour $\frac{3+5}{2\times7}$, saisir \texttt{(3+5)/(2*7)}. L'écriture \texttt{3+5/2*7} désigne autre chose. Pour $\sqrt{13}$, garder la racine dans les étapes suivantes et arrondir seulement le résultat final.
\end{exemple}
Le nombre affiché est parfois arrondi. Une égalité à l'écran ne constitue pas toujours une égalité exacte ; une représentation graphique ne garantit pas non plus l'absence d'autres solutions hors de la fenêtre.
\section{Tableur : cellules et formules}
Une cellule se repère par colonne puis ligne : \texttt{B3}. Une formule commence par \texttt{=}. Les opérations usuelles utilisent \texttt{+}, \texttt{-}, \texttt{*}, \texttt{/}, \texttt{\^}. Les noms de fonctions et les séparateurs dépendent de la langue du logiciel ; les exemples ci-dessous utilisent des noms français et le point-virgule.
\begin{exemple}[Deux tarifs]
\[
\begin{array}{c|lll}
&A&B&C\\\hline
1&\text{Nombre de séances}&\text{Tarif A}&\text{Tarif B}\\
2&0&\texttt{=15+3*A2}&\texttt{=5*A2}\\
3&\texttt{=A2+1}&\texttt{=15+3*A3}&\texttt{=5*A3}
\end{array}
\]
Recopier vers le bas fournit un tableau. Un graphique de type nuage de points représente correctement l'abscisse numérique ; une courbe par catégories peut espacer également des abscisses pourtant inégales.
\end{exemple}
\section{Références relatives et absolues}
En recopiant \texttt{=B2/A2} d'une ligne vers le bas, les références deviennent \texttt{=B3/A3}. Ce sont des références relatives.
Pour fixer une cellule, placer des dollars : \texttt{\$B\$8}. La colonne et la ligne ne changent plus à la recopie.
\begin{exemple}[Fréquences]
Si les effectifs sont dans \texttt{B2:B6}, écrire \texttt{=SOMME(B2:B6)} en \texttt{B8}. En \texttt{C2}, écrire \texttt{=B2/\$B\$8} et recopier jusqu'à \texttt{C6}. Chaque effectif change de ligne, mais le dénominateur reste le total.
\end{exemple}
Une référence mixte, comme \texttt{\$B2}, fixe seulement la colonne. Ce complément aide dans une recopie à deux dimensions.
\begin{center}
\begin{tabularx}{\textwidth}{lX}
\toprule But & Exemple de formule\\\midrule
Somme & \texttt{=SOMME(B2:B6)}\\
Moyenne de données brutes & \texttt{=MOYENNE(A2:A21)}\\
Médiane & \texttt{=MEDIANE(A2:A21)}\\
Étendue & \texttt{=MAX(A2:A21)-MIN(A2:A21)}\\
Moyenne avec effectifs & \texttt{=SOMMEPROD(A2:A6;B2:B6)/SOMME(B2:B6)}\\
Tester une condition & \texttt{=SI(B2<C2;"A";"B ou égalité")}\\
Dé entier simulé & \texttt{=ALEA.ENTRE.BORNES(1;6)}\\\bottomrule
\end{tabularx}
\end{center}
Les fonctions aléatoires peuvent changer à chaque recalcul : pour conserver un résultat expérimental, copier puis coller les valeurs.
\section{Géométrie dynamique : construire des dépendances}
\begin{methode}[Explorer Thalès]
Créer trois points non alignés $A,B,C$ ; créer $M$ \emph{sur} $[AB]$ ; tracer la parallèle à $(BC)$ passant par $M$ ; définir $N$ comme son intersection avec $(AC)$ ; afficher $AM/AB$ et $AN/AC$ ; déplacer $M$ et les sommets.
\end{methode}
Les rapports restent égaux numériquement parce que la construction respecte le parallélisme. Placer $M$ à l'œil près de $[AB]$ ou dessiner une droite « presque parallèle » ne garantit rien.
\begin{attention}
Une construction dynamique correcte préserve les hypothèses lors des déplacements. Des mesures répétées suggèrent une propriété mais ne la démontrent pas : la preuve de Thalès reste celle du chapitre 14.
\end{attention}
Dans un logiciel 3D, construire un cône, un plan parallèle à sa base puis leur intersection permet d'observer les sections. Changer la hauteur du plan et comparer les rayons aide à préparer le raisonnement de proportionnalité.
\section{Problème résolu : réservoir et remplissage}
Un réservoir cylindrique de rayon intérieur $0{,}60\unite m$ et de hauteur $1{,}50\unite m$ est vide. Une pompe fournit un débit constant de $18\unite{L/min}$. On souhaite remplir $80\,\%$ de sa capacité.
\begin{enumerate}
\item \textbf{Capacité.} $V=\pi(0{,}60)^2(1{,}50)=0{,}54\pi\unite{m^3}=540\pi\unite L$.
\item \textbf{Volume visé.} $0{,}8V=432\pi\unite L$.
\item \textbf{Durée.} $t=432\pi/18=24\pi\unite{min}\approx75{,}40\unite{min}$, soit environ $1\unite h\,15\unite{min}\,24\unite s$.
\item \textbf{Fonction.} Avant saturation, le volume ajouté après $t$ minutes est $q(t)=18t$. Cette formule ne décrit pas un volume stocké au-delà de la capacité.
\end{enumerate}
On a utilisé volume, conversion, pourcentage, débit, fonction et domaine de validité.
\section{Problème résolu : une hauteur inaccessible}
Un arbre projette une ombre de $9\unite m$. Un bâton vertical de $1{,}20\unite m$ projette simultanément une ombre de $1{,}80\unite m$, sur le même terrain horizontal. On assimile les rayons lumineux à des droites parallèles.

Les deux triangles ombre--hauteur sont rectangles et possèdent le même angle avec les rayons : ils sont semblables. Donc
\[\frac h9=\frac{1{,}20}{1{,}80}=\frac23,\qquad h=6\unite m.\]
Les hypothèses de verticalité, de terrain horizontal et d'éclairage simultané sont essentielles. Le quotient des deux ombres est un rapport de longueurs ; on ne met pas au carré pour trouver une hauteur.
\section{Problème résolu : optimiser une aire sans dériver}
\approfondissement Pour le rectangle de périmètre $24\unite{cm}$ du chapitre 8 :
\[A(x)=x(12-x)=36-(x-6)^2.\]
Comme un carré est positif ou nul, $A(x)\leq36$ pour $0<x<12$. L'égalité a lieu exactement pour $x=6$. Le rectangle de plus grande aire est donc le carré de côté $6\unite{cm}$.
Cette preuve combine calcul littéral et signe d'un carré ; un graphique peut suggérer le résultat, mais l'identité le démontre.
\section{Une méthode de contrôle universelle}
\begin{enumerate}
\item \textbf{Sens :} le résultat répond-il à la question posée ?
\item \textbf{Conditions :} les hypothèses du théorème ou du modèle sont-elles vérifiées ?
\item \textbf{Unités :} chaque égalité concerne-t-elle des grandeurs compatibles ?
\item \textbf{Ordre de grandeur :} le résultat paraît-il plausible ?
\item \textbf{Exactitude :} ai-je distingué $=$ et $\approx$ ?
\item \textbf{Exhaustivité :} ai-je cherché toutes les solutions ou tous les cas demandés ?
\end{enumerate}
\utilite{Ces contrôles s'appliquent à tous les chapitres et continueront à servir au lycée. Un outil numérique devient utile lorsque l'on sait ce qu'il calcule et ce qu'il ne prouve pas.}
\begin{retenir}
Chercher avec un outil, démontrer avec des arguments, puis contrôler avec un autre regard. Le résultat final est une réponse expliquée, pas seulement un nombre affiché.
\end{retenir}
\end{document}
